%! Introducing Statistics: Do the Data We Collected Tell the Truth? — Unit 3, Lesson 3.1 — Homework
% GRADUAL RELEASE: this is the lesson's SCORED practice and the LAST component in the packet.
% Paper homework is the DEFAULT for this course; DeltaMath is an occasional override that
% replaces this page. Started in the last eight minutes
% of the period and finished after two study halls. The AP Practice component before it is extra credit,
% scored 1 to 5.
% THIRD CONTEXT: the notes teach on a skate-park vote and the school-start poll, Guided Practice
% runs on a fitness app, and these items are on restaurant ratings, college plans, and a gym.
% Part A (restaurant ratings): A1 chance or not, A2 who rates and which way it leans, A3 THE
% CRUX restated — "ten times as many ratings, so report it". Part B: B1 the AP-style multiple
% choice with a required justification, B2 a gym survey's bias in context. 10 points.
% \answerspace{H}{} reserves exactly H; each \gradedblank inside it is BYTE-IDENTICAL in the key,
% where key mode prints its display. Keep every \boxguard identical in the blank and the key.
\documentclass[12pt]{article}
\usepackage{apstats-article}
\usepackage{apstats-boxes}
\usepackage{grading}
\gradingassignment{apstats-u03-l01-hw}
\gradingcourse{advanced_statistics}
\newcommand{\answerspace}[2]{\par\nopagebreak\noindent\begin{minipage}[t][#1][t]{\linewidth}%
  \color{keyred}\bfseries #2\end{minipage}\par}

\begin{document}

\pageheader{Unit 3, Lesson 3.1}{Homework}
{\small\itshape Write each final answer in its box. Only what is inside a box is graded.}\par
% NAMESTRIP: no \namedateperiod here — the name row lives on the cover only.

\vspace{0.04in}

% Authored BYTE-IDENTICALLY in homework/ and homework_key/.
\boxguard[8]
\begin{remindbox}
\small
\textbf{This is your graded homework.} It is scored 1 to 5 and due the first class after two study
halls. Start it now; ask about anything that stalls you before you leave.
\end{remindbox}

\vspace{0.04in}

\boxguard[14]
\begin{notesbox}{Part A --- Restaurant Ratings}
\interlinepenalty=10000 % keep each prompt whole when the box breaks
\small
A restaurant chain wants to know what percent of its customers rate their last meal
``excellent.'' \textbf{Method 1:} a QR code on every receipt invites customers to rate their
meal online; 3{,}200 do, and \textbf{41\%} choose ``excellent.'' \textbf{Method 2:} the chain
picks 300 orders at random from its list of all last month's orders and calls those customers;
\textbf{63\%} say ``excellent.''
\begin{enumerate}[label=\textbf{A\arabic*.}, leftmargin=2.2em, itemsep=2pt]
  \item Which method used chance to decide who answered? In the other method, who decided?
        \answerspace{2.1in}{\normalcolor\mdseries\hspace*{10.5pt}\gradedblank[id=hw1, type=rubric, points=2, display={Method 2: the 300 orders were picked at random. In Method 1 the customers decided, by choosing to scan the code.}]{\linewidth-10.5pt}{0.95cm}}
  \item Who is most likely to scan the code and rate a meal? Explain why Method 1's 41\% is
        likely too low.
        \answerspace{2.1in}{\normalcolor\mdseries\hspace*{10.5pt}\gradedblank[id=hw2, type=rubric, points=2, display={Customers with strong feelings, especially an unhappy one, are most motivated to rate, so they are over-represented and pull the ``excellent'' share down.}]{\linewidth-10.5pt}{0.95cm}}
  \item A manager says, ``Method 1 has more than ten times as many ratings as Method 2, so it is
        the one to report.'' Respond.
        \answerspace{2.1in}{\normalcolor\mdseries\hspace*{10.5pt}\gradedblank[id=hw3, type=rubric, points=2, display={Disagree. 3,200 volunteers are still a biased, non-chance sample; more of them does not remove the bias. Report Method 2's 63\%.}]{\linewidth-10.5pt}{0.95cm}}
\end{enumerate}
\end{notesbox}

\vspace{0.04in}

\boxguard[30]
\begin{notesbox}{Part B --- Chance or Not?}
\interlinepenalty=10000 % keep each prompt whole when the box breaks
\small
\begin{enumerate}[label=\textbf{B\arabic*.}, leftmargin=2.2em, itemsep=2pt]
  \item \textbf{AP-style.} A state wants to estimate the percent of its high school seniors who
        plan to attend college. Which method relies on chance to choose who is surveyed? Write
        its letter in the Letter box, then justify it in a sentence.
        \begin{enumerate}[label=\textbf{(\Alph*)}, leftmargin=2.2em, itemsep=0pt]
          \item Survey the seniors who attend a statewide college fair.
          \item Let each school counselor pick ten seniors to survey.
          \item Post the survey on the state education department's website.
          \item Randomly select 600 seniors from the state's list of all enrolled seniors.
          \item Survey every senior enrolled in an AP class.
                \hfill\gradedright{\textit{Letter:}~\gradedblank[id=hw4, type=choice, points=1, answer=D]{0.45in}{0.24in}}\rule[-0.2in]{0pt}{0.48in}
        \end{enumerate}
        \answerspace{2.1in}{\normalcolor\mdseries\hspace*{10.5pt}\gradedblank[id=hw4j, type=rubric, points=1, display={(D): a random process picks the 600 from the list of all seniors. In the others, people decide.}]{\linewidth-10.5pt}{0.7cm}}
  \item A gym owner wants to know whether members want the gym to stay open later at night. She
        surveys the members who are working out at 6 a.m. Name a group this survey
        over-represents and a group it likely misses, and say which way the result probably
        leans.
        \answerspace{2.1in}{\normalcolor\mdseries\hspace*{10.5pt}\gradedblank[id=hw5, type=rubric, points=2, display={Over-represents early-morning exercisers; misses members who work out in the evening. Support for later hours is likely too low.}]{\linewidth-10.5pt}{0.95cm}}
\end{enumerate}
\end{notesbox}

\vspace{0.04in}

% The next-lesson preview closes the packet, so it lives here rather than in AP Practice.
% A one-line, unbreakable box: guard it by its own height so it moves whole.
\boxguard[3]
\begin{spiralbox}
\small
\textbf{Next: Introduction to Planning a Study} --- populations and samples, a census, and the
difference between an observational study and an experiment.
\end{spiralbox}

\end{document}
